\begin{pastexam}[label={px:ccu-108-swa-05}]{中正資工所甲組 108 學年度 第 6 題}
（問答、證明）

(5\%) Please answer the following questions.
\begin{enumerate}[label=(\alph*)]
\item (2\%) What is the maximum number of tree nodes in a $k$-ary tree of height $h$?
\item (3\%) Prove your answer in (a).
\end{enumerate}
\end{pastexam}

\begin{pastsoln}
答案：(a) 高度數節點（根在第 $1$ 層，資結書第~7.1~節的定義）時，$k \ge 2$ 最多
\[
1 + k + k^2 + \cdots + k^{h-1} = \frac{k^h - 1}{k - 1}
\]
個節點；$k = 1$ 時是 $h$ 個。 題目沒有定義高度。 如果高度數的是邊數（只有根的樹高度 $0$），
同樣的樹有 $h + 1$ 層，答案是 $\dfrac{k^{h+1} - 1}{k - 1}$。 作答時先寫一句用的是哪一種定義。

$k$-ary tree 指每個節點最多 $k$ 個子節點的樹。 $k = 2$、$h = 10$ 時是 $2^{10} - 1 = 1023$，就是
二元樹的資結書定理~7.1。

\textbf{(b)} 證明分兩步：先證每一層的上界，再加起來，最後說明上界達得到。

\textbf{第 $i$ 層最多 $k^{i-1}$ 個節點}，對 $i$ 做數學歸納（資結書第~A.4~節）。
\begin{itemize}
\item 基底 $i = 1$：第 $1$ 層只有根，$1 = k^0$ 個。
\item 歸納步驟：設第 $i$ 層最多 $k^{i-1}$ 個節點。 第 $i + 1$ 層的每個節點都是第 $i$ 層某個節點的
子節點，而每個節點最多 $k$ 個子節點，所以第 $i + 1$ 層最多 $k \cdot k^{i-1} = k^i$ 個節點。
\end{itemize}

\textbf{加起來}。 高度 $h$ 的樹恰好有第 $1$ 到第 $h$ 層，節點總數
\[
n \le \sum_{i=1}^{h} k^{i-1} = 1 + k + \cdots + k^{h-1}.
\]
$k \ge 2$ 時，令 $S = 1 + k + \cdots + k^{h-1}$，則 $kS = k + k^2 + \cdots + k^h$，相減得
$(k - 1)S = k^h - 1$，$S = (k^h - 1)/(k - 1)$。 $k = 1$ 時每一項都是 $1$，$S = h$。

\textbf{上界達得到}。 第 $1$ 到第 $h - 1$ 層的每個節點都恰好有 $k$ 個子節點、第 $h$ 層都是葉節點的
樹（完滿 $k$ 元樹），每一層都恰好 $k^{i-1}$ 個節點，總數正好是 $S$。 所以最大值就是
$(k^h - 1)/(k - 1)$。 例如 $k = 3$、$h = 3$：$1 + 3 + 9 = 13 = (27 - 1)/2$。
\end{pastsoln}
